\documentclass[11pt]{article}

% --- Preamble ---
\usepackage[margin=1in]{geometry}
\usepackage{fancyhdr}
\usepackage{hyperref}
\usepackage{xcolor}
\usepackage{graphicx}
\usepackage{booktabs}
\usepackage{amsmath}
\usepackage{amssymb}
\usepackage{amsthm}
\usepackage{enumitem}
\usepackage{listings}
\usepackage{fontspec}
\setmonofont{JetBrains Mono}
\usepackage{caption}
\usepackage{subcaption}
\usepackage{float}
\usepackage{textcomp}
\usepackage{titling}

% --- Colors ---
\definecolor{linkblue}{RGB}{0, 102, 204}
\definecolor{codebg}{RGB}{245, 245, 245}
\definecolor{codeframe}{RGB}{220, 220, 220}

% --- Convenience commands ---
\newcommand{\menu}[1]{\textsf{#1}}

% --- Hyperref setup ---
\hypersetup{
    colorlinks=true,
    linkcolor=linkblue,
    urlcolor=linkblue,
    citecolor=linkblue,
    pdftitle={Setzer Example Document},
    pdfauthor={Sam Fic},
}

% --- Header / Footer ---
\pagestyle{fancy}
\fancyhf{}
\fancyhead[L]{\small\slshape Setzer Example Document}
\fancyhead[R]{\small\thepage}
\renewcommand{\headrulewidth}{0.4pt}
\renewcommand{\footrulewidth}{0pt}

% --- Theorem environments ---
\theoremstyle{definition}
\newtheorem{definition}{Definition}[section]
\newtheorem{example}[definition]{Example}
\newtheorem{remark}[definition]{Remark}

\theoremstyle{plain}
\newtheorem{theorem}[definition]{Theorem}
\newtheorem{lemma}[definition]{Lemma}
\newtheorem{proposition}[definition]{Proposition}

% --- Listings setup ---
\lstset{
    basicstyle=\ttfamily,
    backgroundcolor=\color{codebg},
    frame=single,
    framerule=0.6pt,
    rulecolor=\color{codeframe},
    numbers=left,
    numberstyle=\tiny\color{gray},
    stepnumber=1,
    breaklines=true,
    captionpos=b,
    xleftmargin=1.5em,
    xrightmargin=1.5em,
    showspaces=false,
    showstringspaces=false,
    showtabs=false,
}

% The `listings` package does not ship a JSON language definition.
% Define a minimal JSON language so the sample configuration highlights correctly.
\lstdefinelanguage{JSON}{
    morekeywords={true,false,null},
    sensitive=false,
    morecomment=[l]{//},
    morestring=[b]",
    morestring=[b]',
    alsoletter={:},
}

% --- Title ---
% Push the title block toward the vertical middle of the cover page.
\pretitle{\begin{center}\vspace*{0.25\textheight}\Huge\bfseries}
\posttitle{\end{center}}
\postdate{\par\end{center}\vspace*{0.2\textheight}}

\title{Example Document: A Tour of \LaTeX{} and Setzer Features}
\author{Sam Fic}
\date{\today}

% --- Document ---
\begin{document}

\maketitle
\thispagestyle{empty}

\clearpage

\tableofcontents
\newpage

% --- Abstract ---
\begin{abstract}
This document is a compact showcase of common \LaTeX{} elements
together with Setzer-specific features and tips.
It is intended as a friendly starting point: you can read the source
to see how each effect is produced, and modify it freely.
Sections cover formatting, mathematics, lists, tables, figures,
code listings, cross-references, and practical guidance for using
the Setzer editor.
\end{abstract}

% ============================================================
\section{Introduction}
\label{sec:intro}

Welcome to \LaTeX{}!
This guide walks through the most useful building blocks of a typical
document.
Each subsection is deliberately short so you can jump straight to the
part you need.

\subsection{What this document contains}
\label{sec:intro:contains}

As summarized in Table~\ref{tab:fruits}, this example covers:

\begin{itemize}[noitemsep]
    \item Text formatting and structure
    \item Mathematics, from inline formulas to multi-line alignments
    \item Various list styles
    \item Professional tables and figures
    \item Code listings
    \item Cross-references, footnotes, and hyperlinks
    \item Setzer editor features and tips
\end{itemize}

% ============================================================
\section{Text Formatting}
\label{sec:text}

\LaTeX{} provides several levels of emphasis.
You can write \emph{emphasized} text, \textbf{bold} text,
\textit{italic} text, \textsc{small caps}, and \underline{underlined}
text.
Combinations such as \textbf{\emph{bold italic}} also work.

\subsection{Lists}

\subsubsection{Itemize}
\label{sec:text:lists:itemize}

A simple bullet list:

\begin{itemize}[noitemsep]
    \item First item
    \item Second item
    \item Third item with a longer description that spans multiple lines
          and demonstrates how itemize wraps text naturally
\end{itemize}

\subsubsection{Enumerate}
\label{sec:text:lists:enumerate}

A numbered list:

\begin{enumerate}[noitemsep]
    \item Step one
    \item Step two
    \item Step three
\end{enumerate}

\subsubsection{Description}
\label{sec:text:lists:description}

A description list is handy for defining terms:

\begin{description}[noitemsep]
    \item[Setzer] A \LaTeX{} editor built with GTK and libadwaita.
    \item[\LaTeX] A high-quality typesetting system for technical documents.
    \item[PDF] Portable Document Format, the usual output of \LaTeX{}.
\end{description}

\subsection{Quotes and Remarks}
\label{sec:text:quotes}

Use \texttt{quote} for short quotations and \texttt{quotation} for
longer ones.
Here is a short quote:

\begin{quote}
    \small
    The \TeX{} typesetting system was designed by Donald Knuth
    and first released in 1978.
\end{quote}

A remark environment is useful for side notes:

\begin{remark}
    You do not need to memorize all \LaTeX{} commands.
    Cheat sheets and editors with autocompletion help a lot.
\end{remark}

% ============================================================
\section{Mathematics}
\label{sec:math}

\LaTeX{} excels at typesetting mathematics.
This section shows inline formulas, displayed equations, and
multi-line alignments.

\subsection{Inline Math}
\label{sec:math:inline}

The quadratic formula is $x = \frac{-b \pm \sqrt{b^2 - 4ac}}{2a}$.
You can also use \(\frac{\partial u}{\partial t}\) inline without
breaking the paragraph flow.

\subsection{Displayed Equations}
\label{sec:math:display}

Einstein's famous relation:

\begin{equation}
    \label{eq:einstein}
    E = mc^2
\end{equation}

A system of linear equations:

\begin{align}
    \label{eq:linear}
    f(x) &= (x+a)(x+b) \nonumber \\
         &= x^2 + (a+b)x + ab
\end{align}

A piecewise definition:

\begin{equation}
    \label{eq:piecewise}
    |x| =
    \begin{cases}
        x,  & \text{if } x \ge 0, \\
        -x, & \text{if } x < 0.
    \end{cases}
\end{equation}

A matrix:

\begin{equation}
    \label{eq:matrix}
    A =
    \begin{pmatrix}
        a_{11} & a_{12} & \cdots & a_{1n} \\
        a_{21} & a_{22} & \cdots & a_{2n} \\
        \vdots & \vdots & \ddots & \vdots \\
        a_{m1} & a_{m2} & \cdots & a_{mn}
    \end{pmatrix}
\end{equation}

\subsection{Theorems and Proofs}
\label{sec:math:theorems}

\begin{theorem}[Pythagoras]
    \label{thm:pythagoras}
    In a right-angled triangle, $a^2 + b^2 = c^2$, where $c$ is the
    hypotenuse.
\end{theorem}

\begin{proof}
    This is the classic Euclidean proof.
    Many algebraic proofs also exist.
\end{proof}

\begin{lemma}
    \label{lem:trichotomy}
    For any real number $x$, exactly one of the following holds:
    $x > 0$, $x = 0$, or $x < 0$.
\end{lemma}

\begin{definition}[Prime Number]
    \label{def:prime}
    A prime number is a natural number greater than $1$ that has no
    positive divisors other than $1$ and itself.
\end{definition}

\begin{example}
    \label{ex:primes}
    The first six prime numbers are $2, 3, 5, 7, 11, 13$.
\end{example}

% ============================================================
\section{Tables}
\label{sec:tables}

Table~\ref{tab:fruits} uses the \texttt{booktabs} package for
professional horizontal rules.

\begin{table}[H]
    \centering
    \caption{Example Table}
    \label{tab:fruits}
    \begin{tabular}{llr}
        \toprule
        \textbf{Item} & \textbf{Description} & \textbf{Price} \\
        \midrule
        Apple  & Fresh fruit      & \$1.00 \\
        Banana & Yellow fruit     & \$0.50 \\
        Cherry & Small red fruit  & \$2.00 \\
        \midrule
        \multicolumn{2}{r}{\textbf{Total:}} & \$3.50 \\
        \bottomrule
    \end{tabular}
\end{table}

A second table with aligned columns:

\begin{table}[H]
    \centering
    \caption{Notation Summary}
    \label{tab:notation}
    \begin{tabular}{p{3.5cm} p{6cm}}
        \toprule
        \textbf{Symbol} & \textbf{Meaning} \\
        \midrule
        $\mathbb{R}$  & The set of real numbers \\
        $\forall$     & Universal quantifier, ``for all'' \\
        $\exists$     & Existential quantifier, ``there exists'' \\
        $\implies$    & Logical implication \\
        $\iff$        & Logical equivalence \\
        $\partial$    & Partial derivative \\
        $\nabla$      & Gradient operator \\
        $\infty$      & Infinity \\
        $\emptyset$   & Empty set \\
        $\approx$     & Approximately equal \\
        $\sim$        & Distributed as \\
        $\propto$     & Proportional to \\
        \bottomrule
    \end{tabular}
\end{table}

% ============================================================
\section{Figures and Captions}
\label{sec:figures}

Figure~\ref{fig:placeholder} shows a placeholder box.
In a real document, replace it with \texttt{\textbackslash
includegraphics}.

\begin{figure}[H]
    \centering
    \fbox{\parbox{0.55\textwidth}{\centering
        Figure placeholder\\
        Insert image with \texttt{\textbackslash includegraphics}
    }}
    \caption{Example Figure}
    \label{fig:placeholder}
\end{figure}

You can also place figures side by side using \texttt{subcaption}:

\begin{figure}[H]
    \centering
    \begin{subfigure}[b]{0.45\textwidth}
        \centering
        \fbox{\parbox{0.8\linewidth}{\centering Sub-figure A}}
        \caption{First sub-figure}
        \label{fig:sub:a}
    \end{subfigure}
    \hfill
    \begin{subfigure}[b]{0.45\textwidth}
        \centering
        \fbox{\parbox{0.8\linewidth}{\centering Sub-figure B}}
        \caption{Second sub-figure}
        \label{fig:sub:b}
    \end{subfigure}
    \caption{Multiple sub-figures}
    \label{fig:subfigures}
\end{figure}

% ============================================================
\section{Code Listings}
\label{sec:code}

The \texttt{listings} package renders source code with syntax-aware
highlighting.

\subsection{Python}
\label{sec:code:python}

\begin{lstlisting}[language=Python, caption={A simple Python function.}]
def fibonacci(n: int) -> int:
    """Return the n-th Fibonacci number."""
    if n < 0:
        raise ValueError("n must be non-negative")
    a, b = 0, 1
    for _ in range(n):
        a, b = b, a + b
    return a

if __name__ == "__main__":
    for i in range(10):
        print(f"F({i}) = {fibonacci(i)}")
\end{lstlisting}

\subsection{JSON}
\label{sec:code:json}

\begin{lstlisting}[language=JSON, caption={A sample configuration file.}]
{
    "name": "setzer-example",
    "version": "1.0.0",
    "dependencies": {
        "gtk": ">=4.0",
        "libadwaita": ">=1.0"
    },
    "settings": {
        "theme": "auto",
        "fontSize": 12,
        "autoSave": true
    }
}
\end{lstlisting}

% ============================================================
\section{Cross-References}
\label{sec:references}

This document is heavily cross-referenced so you can see how labels
and references work together.

\subsection{Links}
\label{sec:references:links}

Equation~(\ref{eq:einstein}) expresses mass--energy equivalence.
Theorem~\ref{thm:pythagoras} appears in Section~\ref{sec:math:theorems}.
Definition~\ref{def:prime} and Example~\ref{ex:primes} are in the same
section.
Table~\ref{tab:fruits} and Table~\ref{tab:notation} are in
Section~\ref{sec:tables}.
Figure~\ref{fig:placeholder} and the sub-figures
(\ref{fig:sub:a}, \ref{fig:sub:b}, \ref{fig:subfigures}) are in
Section~\ref{sec:figures}.

You can also create clickable hyperlinks:
\url{https://www.latex-project.org/} and
\href{https://github.com/}{GitHub}.

\subsection{Footnotes}
\label{sec:references:footnotes}

\LaTeX{} supports footnotes easily.\footnote{This is a footnote.}
You can add as many as you need, and \LaTeX{} will place them at the
bottom of the page.\footnote{Another footnote with a citation:
\cite{knuth1986}.}

% ============================================================
\section{Setzer Features and Tips}
\label{sec:setzer}

This section highlights features that make Setzer a comfortable
environment for writing \LaTeX{} on Linux.
All shortcuts shown are the defaults; they can be customized in
\menu{Preferences}.

\subsection{Editing Comfort}
\label{sec:setzer:editing}

\begin{itemize}[noitemsep]
    \item \textbf{Syntax highlighting} with automatic dark/light theme
          support.
    \item \textbf{Code folding} for \texttt{\textbackslash begin\ldots
          \textbackslash end} blocks; folding state is preserved between
          sessions.
    \item \textbf{Smart brackets}: typing \texttt{\textbraceleft\textbraceright}, (, [, \$ inserts
          the matching closing bracket. Select text first to wrap it.
    \item \textbf{Auto-indent} and line operations:
          \menu{Alt+Up / Alt+Down} moves lines,
          \menu{Ctrl+Shift+K} deletes a line,
          \menu{Alt+Shift+D} duplicates it.
\end{itemize}

\subsection{Completion and Snippets}
\label{sec:setzer:completion}

Setzer offers context-aware completion:
type \texttt{\textbackslash} to open the command list, or
\texttt{Ctrl+Space} to trigger it manually.
In math mode, only math-relevant commands are shown.
The \textbf{LaTeXDB} index scans your open documents for
\texttt{\textbackslash label}, \texttt{\textbackslash ref},
\texttt{\textbackslash cite}, and \texttt{\textbackslash bibitem}
commands, so cross-references and citations are completed dynamically.

\subsection{Shortcuts Bar}
\label{sec:setzer:shortcuts}

The sidebar shortcuts bar gives quick access to common insertions
without memorizing every command:
\begin{itemize}[noitemsep]
    \item \textbf{Text menu}: font styles, sections, lists, quotes,
          cross-references.
    \item \textbf{Math menu}: environments, operators, Greek letters,
          accents, fractions.
    \item \textbf{Document menu}: \texttt{\textbackslash documentclass},
          packages, title metadata, TOC.
    \item \textbf{Bibliography menu}: \texttt{\textbackslash bibliography},
          \texttt{\textbackslash bibliographystyle}, \texttt{\textbackslash nocite}.
\end{itemize}

\subsection{Build and Preview}
\label{sec:setzer:build}

Setzer supports PDFLaTeX, XeLaTeX, LuaLaTeX, Tectonic, and \texttt{latexmk}.
Key shortcuts:
\begin{itemize}[noitemsep]
    \item \textbf{F5}: save and build.
    \item \textbf{F6}: build only.
    \item \textbf{F7}: SyncTeX forward search (preview $\to$ source).
    \item \textbf{Ctrl+click} in the preview jumps back to the source.
\end{itemize}

The build log parser categorizes errors, warnings, and bad boxes;
click a log entry to jump to the offending line.

\subsection{AI Fix}
\label{sec:setzer:ai-fix}

Setzer can connect to external AI agent CLIs to suggest fixes for
\LaTeX{} build errors directly from the build log.
When an error is selected, Setzer assembles a prompt containing the
error message and a snippet of surrounding source code, then hands it
off to an agent such as \texttt{opencode}, \texttt{claude},
\texttt{codex}, \texttt{gemini}, or \texttt{aider}.
The agent runs in a real terminal window so you can watch every step;
Setzer saves the document first, then reloads it automatically after
the agent finishes.
For trusted project directories, the hand-off happens without any
extra confirmation dialogs.

\subsection{Session Management}
\label{sec:setzer:session}

Setzer can restore your previous session with
\menu{Ctrl+Shift+J}.
Auto-save writes unsaved changes to a temporary location every few
seconds, and a recovery dialog appears after a crash.
You can also reopen a recently closed document with
\menu{Ctrl+Shift+T}.

\subsection{AI-Driven Workflow}
\label{sec:setzer:ai-workflow}

For an AI-driven writing and compilation workflow, combine
\textbf{AI Fix} with two editor preferences:
\begin{enumerate}[noitemsep]
    \item Open \menu{Preferences} and go to the
          \menu{Editor $\rightarrow$ External Changes} group.
          Enable \menu{Automatically reload changed files} so that
          Setzer picks up edits made by external AI agents without
          prompting.
    \item Then enable \menu{Auto Build} in the
          \menu{Build System} preferences.
          The document will be saved and rebuilt automatically
          shortly after you stop typing.
\end{enumerate}
With these settings, an external AI agent can edit your \LaTeX{}
source and Setzer will immediately reload and recompile the
document, giving you a smooth AI-assisted authoring loop.

\subsection{Tips for New Users}
\label{sec:setzer:tips}

\begin{enumerate}[noitemsep]
    \item Use the \textbf{document wizard} from the welcome screen to
          generate a preamble with sensible defaults.
    \item Set a \textbf{root document} for multi-file projects so that
          \menu{Build} always compiles the right entry point.
    \item Enable \textbf{auto-build} in preferences to rebuild
          automatically after you stop typing.
    \item Explore the built-in \textbf{LaTeX help panel} (F1) for
          searchable LaTeX2e documentation.
    \item Use \textbf{pinned recent documents} to keep important files
          at the top of the welcome screen list.
\end{enumerate}


\subsection{More Shortcuts}
\label{sec:setzer:shortcuts-more}

A quick reference for common shortcuts:

\begin{itemize}[noitemsep]
    \item \textbf{Undo / Redo}: \menu{Ctrl+Z} and \menu{Ctrl+Shift+Z}.
    \item \textbf{Comment / Uncomment}: \menu{Ctrl+/} toggles line comments.
    \item \textbf{Duplicate line}: \menu{Alt+Shift+D} copies the current line.
    \item \textbf{Jump to line}: open the jump dialog from the menu.
    \item \textbf{Full screen}: \menu{F11} toggles full-screen mode.
    \item \textbf{Zoom editor}: \menu{Ctrl+Plus}, \menu{Ctrl+Minus}, and \menu{Ctrl+0}.
    \item \textbf{Help panel}: \menu{F1} opens the built-in LaTeX help browser.
    \item \textbf{Document structure}: \menu{Ctrl+Shift+B} shows the section outline.
    \item \textbf{Symbols panel}: \menu{F8} opens the math symbols sidebar.
\end{itemize}


% ============================================================
\section{About This Fork}
\label{sec:about}

This example document is maintained as part of a fork of Setzer.

\begin{itemize}[noitemsep]
    \item \textbf{Fork repository}: \url{https://github.com/Sam-Fic/Setzer}
    \item \textbf{Upstream project}: \url{https://github.com/cvfosammmm/setzer}
\end{itemize}

Development and discussion for this fork take place on GitHub.
For the original upstream project, see the link above.

% ============================================================
\section{Conclusion}
\label{sec:conclusion}

This example demonstrates the core elements of a \LaTeX{} document.
From here you can:

\begin{enumerate}[noitemsep]
    \item Open the source and change sections, labels, or captions.
    \item Add your own packages and environments.
    \item Export to PDF and inspect the rendered output.
\end{enumerate}

Happy \TeX ing!

% ============================================================
\bibliographystyle{plain}
\begin{thebibliography}{9}

\bibitem{knuth1986}
Donald~E. Knuth.
\newblock {\em TeX: The Program}.
\newblock Addison-Wesley, Reading, Massachusetts, 1986.

\end{thebibliography}

\end{document}
